% ===================================================================
%  TAULA N° 16 — Los Grands Magazins. — Lo Basar. — Las Joguinas.
%  Traduction 2026 d'apres texto/io, controlee sur texto/fr et
%  texto/ca. Consigne : voir texto/oc/10-taula-01.tex.
%
%  Le tableau d'astronomie (bloc 15-1) porte des renvois a LETTRES,
%  (a) a (f), graves sur le tableau lui-meme. L'ido saute la lettre
%  (d), que le fac-simile grave pourtant : la colonne la rend, comme
%  toutes les autres traductions. Ecart declare dans outils/renvoji.py.
%
%  DERNIER TABLEAU DE LA COLONNE OCCITANE. Bilan des seize : zero
%  inversion, comme en catalan. Les deux colonnes sœurs ont le meme
%  compte, et pour la meme raison — l'ido a pris son ordre des mots
%  aux langues romanes.
% ===================================================================

%%K t16-apar-1 apar
\VUsaut{4.0mm}
\VUtitre{15.0pt}{60}{TAULA N° 16}
\VUsaut{2.0mm}
\VUpk{13.2pt}{40}{Los Grands Magazins. --- Lo Basar.}
\VUpk{13.2pt}{40}{Las Joguinas.}
\VUsaut{3.4mm}

%%K t16-01-1 p
1. --- S'apròcha l'an novèl. Los grands magazins an preparat lors exposicions de
\VUgras{joguinas} \textsuperscript{(1)} e de presents d'an novèl.

%%K t16-01-2 p
Demandèri a mon grand que me menèsse al basar veire aquela bèla exposicion, e
acceptèt.

%%K t16-02-1 p
2. --- D'en primièr nos arrestèrem davant de bèlas panoplias d'armas amb un
\VUgras{fusilh}, una \VUgras{casqueta}, un \VUgras{sabre}, un \VUgras{ceinturon} e
un parelh d'\VUgras{espauletas}, o ben un \VUgras{casco}, una \VUgras{cuirassa}
\textsuperscript{(3)} e un \VUgras{pistolet} \textsuperscript{(4)}. Tot aquò
m'interessava fòrça mai que las \VUgras{lanternas magicas}
\textsuperscript{(2)}.

%%K t16-tit-1 sub
\VUsaut{3.0mm}
\VUpk{10.2pt}{40}{Bèles espectacles.}
\VUsaut{2.6mm}

%%K t16-03-1 p
3. --- Dins la meteissa seccion i aviá una \VUgras{sentinèla}
\textsuperscript{(5)} dins sa \VUgras{garita}; semblava montar la garda sus tot un
parc de fèras de carton. Quantes animals i aviá! Un \VUgras{papagai}
\textsuperscript{(6)} sus son baston, un \VUgras{rinocèros} \textsuperscript{(7)}
amb una bana al nas, un \VUgras{ren} \textsuperscript{(8)}, un
\VUgras{estruci} \textsuperscript{(9)}, una \VUgras{mona} \textsuperscript{(10)}
que copava una nose, un \VUgras{rainal} \textsuperscript{(11)} qu'emportava una
galina, un \VUgras{cocodril} \textsuperscript{(12)} amb sas enòrmas maissas
dobèrtas, un \VUgras{camèl} \textsuperscript{(13)}, un elegant \VUgras{levrièr}
\textsuperscript{(14)}, una \VUgras{pantèra} \textsuperscript{(15)}, e puèi doas
bèstias que ieu coneissiái pas, que son, segon me dison, una \VUgras{moisseta}
\textsuperscript{(16)} e una \VUgras{iena} \textsuperscript{(17)}. I aviá tanben un
\VUgras{leopard} \textsuperscript{(18)} e un \VUgras{lop} \textsuperscript{(21)};
fòrça prèp i aviá de polidas \VUgras{ratas} \textsuperscript{(19)} e de
\VUgras{ratons} \textsuperscript{(20)} pichonets de cauchoc.

%%K t16-03-1 p suite
Enfin, un \VUgras{ipopotam} \textsuperscript{(22)} e un \VUgras{dromadari}
\textsuperscript{(23)} bossut completavan la colleccion. I seriái demorat d'oras
entièras s'i aviá pas agut, fòrça prèp, un esplendid campament plen de
\VUgras{soldadets de plomb} \textsuperscript{(29)} que me cridèt l'atencion.

%%K t16-tit-2 sub
\VUsaut{4.0mm}
\VUpk{10.2pt}{40}{Soldats e guèrra.}
\VUsaut{2.8mm}

%%K t16-04-1 p
4. --- Mon grand, qu'es un \VUgras{invalid de guèrra} \textsuperscript{(24)} e
deguèt daissar l'armada a causa d'una \VUgras{nafradura} que li valguèt la
\VUgras{medalha militara}, adòra me contar las \VUgras{campanhas} ont participèt.
L'encantèt de m'explicar los diferents movements de la pichona armada en
miniatura. Un grop de soldats de vertat que visitavan lo basar --- un
\VUgras{enginhaire militar} \textsuperscript{(25)}, un \VUgras{caçaire alpin}
\textsuperscript{(26)} amb sa \VUgras{bereta}, un \VUgras{sergent major}
\textsuperscript{(27)} d'infantariá e un \VUgras{sergent} \textsuperscript{(28)}
d'artilhariá amb un \VUgras{galon} d'aur sus la margue --- s'arrestèt contra
nosautres per escotar las explicacions.

%%K t16-05-1 p
5. --- Mon grand, fòrça fièr d'aver un public tan distinguit, improvisèt tot un
plan de batalha.

%%K t16-05-2 p
La plaça fòrta, amb sas \VUgras{muralhas} e sos \VUgras{fossats}, èra assetjada
per l'\VUgras{armada enemiga}, que, aprèp un long \VUgras{bombardament},
s'aprestava a la prene d'\VUgras{assaut}. Mas los \VUgras{assetjats}, butats per
la fam, avián ensajat una \VUgras{sortida} contra lo \VUgras{campament} dels
\VUgras{assetjaires}. Repossat l'\VUgras{ataca} dins un \VUgras{combat}
aferonit a l'entorn de las \VUgras{tendas}, los assetjaires \VUgras{vencedors}
lancèron lors \VUgras{esquadrons} a la \VUgras{persecucion} dels assaltants
\VUgras{vencuts}. Aqueles \VUgras{se retirèron}, en daissant lo \VUgras{camp de
batalha} semenat de \VUgras{mòrts} e de \VUgras{nafrats}. Los
\VUgras{brancardièrs} menèron los nafrats a l'\VUgras{ambulància} per que los
suenhèsse lo \VUgras{cirurgian}.

%%K t16-05-3 p
Malgrat l'eroïca resisténcia de qualques \VUgras{batalhons} que s'èran formats en
\VUgras{carrat}, la \VUgras{desfacha} foguèt completa; la rèsta de l'armada
\VUgras{fugiguèt}, en daissant darrièr fòrça \VUgras{presonièrs}. L'enemic èra a
mand de coronar sa \VUgras{victòria} en prenent la plaça.

%%K t16-05-3 p suite
Benlèu que seriá aquò la fin de la campanha, e aprèp un \VUgras{armistici} se
poiriá signar un \VUgras{tractat de patz}.

%%K t16-tit-3 sub
\VUsaut{4.4mm}
\VUpk{10.2pt}{40}{Los presents d'an novèl.}
\VUsaut{2.8mm}

%%K t16-06-1 p
6. --- Los joves soldats, que risián en escotar lo pintoresc raconte del veteran,
li faguèron qualques questions de mai. Quant a ieu, faguèri lo torn del comptador
per agachar una bèla \VUgras{arbalèsta} \textsuperscript{(32)} e un \VUgras{arc}
\textsuperscript{(33)} amb sa \VUgras{sageta} \textsuperscript{(31)}. Aquí trobèri
una autra colleccion d'animals: dos \VUgras{aucèls cantaires}
\textsuperscript{(34)}, un \VUgras{rossinhòl} e una \VUgras{fauveta} de mecanica
que cantavan a plen gargalhòl, e una seguida de bèstias estranhas --- una
\VUgras{ratapenada} \textsuperscript{(35)}, un \VUgras{boà}
\textsuperscript{(36)}, una \VUgras{tartuga} \textsuperscript{(37)}, una
\VUgras{fòca} \textsuperscript{(38)}, una \VUgras{balena} \textsuperscript{(39)} e
un \VUgras{tauron} \textsuperscript{(40)} faches de pèl daurada.

%%K t16-07-1 p
7. --- M'arrestèri pas gaire a los agachar, qu'aviái vist çò que somiavi: un
magnific \VUgras{teatre de marionetas} \textsuperscript{(41)} amb d'esplendids
\VUgras{decòrs} e un deliciós \VUgras{ridèu} roge. Sus l'\VUgras{empont}, dos
actors semblavan representar un \VUgras{ròtle} dins una \VUgras{pèça}
interessantissima, que los pichons \VUgras{espectators} de carton de las lòjas, o
assetats dins los \VUgras{fautuèlhs} e dins lo \VUgras{parterre} darrièr lo
\VUgras{director d'orquèstra}, aplaudissián amb entosiasme.

%%K t16-07-2 p
Per malastre, aquel teatre èra caríssim.

%%K t16-08-1 p
8. --- En mai, èra malaisit de causir entre las joguinas, que las vitrinas èran
plenas de joguinas de tota mena: un \VUgras{Arlequin} \textsuperscript{(42)}, un
\VUgras{Polichinèla} \textsuperscript{(43)}, un \VUgras{Pierrot}
\textsuperscript{(44)}, de \VUgras{cavècas} \textsuperscript{(45)} que batián de
las alas, de \VUgras{mèrles} \textsuperscript{(46)} que siblavan, de
\VUgras{canichs} que japavan, un \VUgras{gramofòn} \textsuperscript{(47)} e de
\VUgras{casse-cap}.

%%K t16-noto-1 noto
\VUnotes{143.2mm}{%
\textit{(*) Vejatz la taula n° 6.}%
}

%%K t16-08-1 p suite
Una d'aquelas joguinas me divertiguèt fòrça: representava un \VUgras{combat naval}
\textsuperscript{(48)} ont un \VUgras{naviri} es atacat per de \VUgras{pirates}
davant una còsta plantada de \VUgras{cocotièrs}.

%%K t16-09-1 p
9. --- Poguèri agachar totas aquelas meravilhas amb tota calma, que i aviá pauc de
mond dins los magazins --- tot bèl just un ponhat de curioses, un \VUgras{dragon}
\textsuperscript{(49)} e un \VUgras{recobraire} \textsuperscript{(50)} que
presentava una \VUgras{letra de cambi} a la \VUgras{caissa}
\textsuperscript{(51)} principala.

%%K t16-10-1 p
10. --- Demandèri a mon grand de qué servissián los libres qu'i aviá suls
estatgièrs darrièr lo \VUgras{caissièr}; me diguèt qu'èran los \VUgras{libres de
comptes}.

%%K t16-tit-4 sub
\VUsaut{3.6mm}
\VUpk{10.2pt}{40}{En badar per tota la botiga.}
\VUsaut{2.8mm}

%%K t16-11-1 p
11. --- En sortir de la seccion de joguinas, anèrem a la selariá getar una uèlhada
a las \VUgras{selas} \textsuperscript{(53)}, als \VUgras{estriups}
\textsuperscript{(54)}, a las \VUgras{bridas} \textsuperscript{(55)}, als
\VUgras{mòrs} \textsuperscript{(56)} e a las \VUgras{cravachas}. Mentre que lo
\VUgras{cap de seccion} \textsuperscript{(57)} donava a mon grand qualques
entresenhas, anèri agachar l'exposicion de \VUgras{porcelana} e de
\VUgras{cristalariá} \textsuperscript{(58)}, ont i aviá de bèlas
\VUgras{ensaladièras} e de finas \VUgras{tetièras} \textsuperscript{(59)}.

%%K t16-12-1 p
12. --- Per pujar al primièr estatge, passèrem darrièr la vitrina dels
\VUgras{corsets} \textsuperscript{(60)}, contra la merçariá, ont èran expausats de
\VUgras{calçons} \textsuperscript{(63)} fòrça bèles. A prèp, una
\VUgras{vendeirèla} \textsuperscript{(61)} ajudava una \VUgras{crompaira}
\textsuperscript{(62)} a causir de bèlas \VUgras{telas} \textsuperscript{(64)} de
seda.

%%K t16-12-2 p
En pujar l'escalièr, notèri que los magazins èran ornats de \VUgras{lanternas}
\textsuperscript{(65)} japonesas, d'\VUgras{ombrèlas} \textsuperscript{(66)} e
d'enòrmas \VUgras{palmièras} \textsuperscript{(70)}.

%%K t16-13-1 p
13. --- Al primièr estatge i aviá pas gaire causa de veire: pas que de mòbles (*),
de \VUgras{còfres-fòrts} \textsuperscript{(67)}, de \VUgras{comòdas}
\textsuperscript{(68)} e de \VUgras{cochets d'enfant} \textsuperscript{(69)}. Mas
la vista generala dels magazins èra fòrça \VUgras{divertenta}.

%%K t16-14-1 p
14. --- Jos nosautres vesiái los instruments de musica: de \VUgras{arpas}
\textsuperscript{(71)}, de \VUgras{guitarras} \textsuperscript{(72)}, de
\VUgras{mandolinas} \textsuperscript{(73)}, de \VUgras{cornets a pistons}
\textsuperscript{(74)}, de \VUgras{clarinetas} \textsuperscript{(75)}, de violons
amb lors \VUgras{arquets} \textsuperscript{(76)} e quitament una \VUgras{grossa
caissa} \textsuperscript{(77)}. Un pauc mai luènh, al comptador de papetariá, un
\VUgras{estudiant} \textsuperscript{(78)} marchandava amb lo \VUgras{comés}
\textsuperscript{(79)} una cartabèla de documents. Contra eles destriavi un bèl
\VUgras{abecedari} \textsuperscript{(80)}.

%%K t16-14-2 p
Al mitan dels magazins s'auçava un naut \VUgras{arbre de Nadal}
\textsuperscript{(81)} cargat de candelas, de tissas e de joguinas de tota mena:
de \VUgras{cavals de fusta} \textsuperscript{(82)}, de \VUgras{pipautas}
\textsuperscript{(83)}, etc.

%%K t16-14-3 p
La \VUgras{parfumariá} \textsuperscript{(84)}, amb sos \VUgras{dentifricis} e sos
divèrses \VUgras{parfums}, escampava una odor fòrça agradiva que pujava fins a
nosautres.

%%K t16-tit-5 sub
\VUsaut{3.4mm}
\VUpk{10.2pt}{40}{Crompam quicòm.}
\VUsaut{2.8mm}

%%K t16-15-1 p
15. --- Coma mon grand començava de trobar lo temps long, tornèrem davalar per lo
grand escalièr. S'arrestèt un moment davant un \VUgras{quadre d'astronomia}
\textsuperscript{(85)} que representava, çò que me diguèt, de \VUgras{cometas}
\textsuperscript{(a)}, d'\VUgras{eclipsis de luna e de solelh}
\textsuperscript{(b)}, una ròsa dels vents amb los \VUgras{quatre punts cardinals}
\textsuperscript{(c)} \VUgras{(nòrd, sud, èst e oèst)}, una mapa dels dos
\VUgras{emisfèris} \textsuperscript{(d)}, los \VUgras{planetas}
\textsuperscript{(e)} e lo \VUgras{sistèma solar} \textsuperscript{(f)}.
Comprenguèri res d'aquò, e m'interessèron fòrça mai la liurèa galonada d'un
\VUgras{liuraire} \textsuperscript{(86)} que passava, e los polits
\VUgras{ventalhs} \textsuperscript{(87)} qu'aviái al costat, contra los
\VUgras{portamonedas} \textsuperscript{(88)} e los \VUgras{portafuèlhas}
\textsuperscript{(89)}.

%%K t16-16-1 p
16. --- Admirèri tanben sul comptador de \VUgras{plumas} \textsuperscript{(90)} e
de \VUgras{fivèlas de ceinturon} \textsuperscript{(91)}, e las polidas
\VUgras{flors artificialas} \textsuperscript{(94)} graciosament dispausadas dins
de panièrs.

%%K t16-16-1 p suite
Las \VUgras{violetas}, los \VUgras{violièrs}, los \VUgras{jacints}
\textsuperscript{(95)}, los \VUgras{geraniums} \textsuperscript{(96)}, los
\VUgras{liris} \textsuperscript{(97)} e las \VUgras{capucinas}
\textsuperscript{(98)} èran fòrça plan imitats.

%%K t16-tit-6 sub
\VUsaut{4.4mm}
\VUpk{10.2pt}{40}{Lo present del grand.}
\VUsaut{2.8mm}

%%K t16-17-1 p
17. --- Demandèri a mon grand que me crompèsse una de las \VUgras{bròssas de
dents} \textsuperscript{(92)} d'una caissa qu'i aviá contra las \VUgras{espinglas
de cap} \textsuperscript{(93)}. Acceptèt, e anèri ieu meteis pagar ma crompa a la
caissa, al costat de la quala fasián un grand \VUgras{paquet}
\textsuperscript{(100)} per un \VUgras{client} \textsuperscript{(99)}.

%%K t16-18-1 p
18. --- Tot d'un còp i aguèt un leugièr bolegadís entre lo mond qu'èra contra los
\VUgras{bronzes} \textsuperscript{(101)} artistics. Un \VUgras{panaire}
\textsuperscript{(102)} veniá d'èsser suspres en flagrant delicte per un
\VUgras{inspector dels magazins} \textsuperscript{(103)} quand ensajava de raubar
a qualqu'un. Se mandèt quèrre un garda per l'\VUgras{arrestar}, e al moment ont
sortissiam, emportavan lo colpable a la comissariá.
\VUsaut{6.0mm}
\VUfilet{22mm}
